\chapter{Géométrie}
% \section*{Exemple introductif}
% % Placez ici un exemple motivant illustrant la technique.
%
% \section*{Résumé de la méthode}
% % Résumez ici les idées clés de la méthode avec des repères pour l'élève.
%
% \section*{Exemple applicatif}
% % Un exercice résolu pas à pas montrant l'application de la méthode.
%
% \section*{Exercices guidés}
% % Exercices progressifs
%
%
% \section*{Problèmes complexes}
% % Sélection de problèmes difficiles pour entraîner la maîtrise du chapitre.

\section*{Notions utiles}

Dans tout le chapitre, les points sont dans le plan ou dans l'espace. $AB$ ou $|AB|$ désigne la longueur du segment $[AB]$.

\textbf{Vecteur} : Un vecteur est défini par une direction, un sens et une norme (sa longueur). On note $\vect{AB}$ le vecteur :
    \begin{itemize}
        \item de direction $(AB)$,
        \item de sens allant du point $A$ au point $B$,
        \item de longueur $AB$.
    \end{itemize}
Relation de Chasles : $\vect{AB} + \vect{BC} = \vect{AC}$. Si l'on fixe une origine $O$, on identifie un point $X$ au vecteur $\vect{OX}$ et on le note simplement $X$ : alors $\vect{AB} = B - A$.

\textbf{Produit scalaire} : $\vec{u}\cdot\vec{v}=\|\vec{u}\|\,\|\vec{v}\|\cos\theta$, où $\theta$ est l'angle entre $\vec{u}$ et $\vec{v}$. On a $\vec{u}\cdot\vec{u}=\|\vec{u}\|^2$, et $\vec{u}\cdot\vec{v}=0$ si et seulement si $\vec{u}$ et $\vec{v}$ sont orthogonaux (ou l'un d'eux est nul). Il est bilinéaire, d'où $\|\vec{u}\pm\vec{v}\|^2=\|\vec{u}\|^2\pm2\,\vec{u}\cdot\vec{v}+\|\vec{v}\|^2$.

\textbf{Milieu d'un segment} : Le milieu du segment $[AB]$ est le point $M$ tel que $\vect{AM} = \vect{MB}$. Ses coordonnées sont $ \left( \dfrac{x_A+x_B}{2} ; \dfrac{y_A+y_B}{2} ; \ldots \right)$ ; avec l'identification ci-dessus, $M=\dfrac{A+B}{2}$.

\textbf{Centroïde (centre de gravité)} : Le centroïde de $n$ points $A_1, A_2, \ldots, A_n$ est l'unique point $S$ tel que
\[
\sum_{i=1}^{n} \vect{SA_i} = \vec{0},
\quad\text{c'est-à-dire}\quad S=\frac{A_1+\cdots+A_n}{n}.
\]
Pour un triangle, c'est le point d'intersection des médianes (les droites joignant un sommet au milieu du côté opposé).

\textbf{Parallélogramme} : Un quadrilatère $ABCD$ dont les côtés opposés sont parallèles. Il est équivalent de demander $\vect{AB}=\vect{DC}$, ou que les diagonales $[AC]$ et $[BD]$ aient le même milieu, c'est-à-dire $A+C=B+D$. Dans l'espace, on dit « parallélogramme » aussi pour un parallélogramme aplati (points alignés).

\textbf{Médianes d'un quadrilatère} : dans un quadrilatère $ABCD$, ce sont les deux segments qui joignent les milieux de côtés opposés : celui de $[AB]$ à celui de $[CD]$, et celui de $[BC]$ à celui de $[DA]$.

\textbf{Orthocentre et cercle circonscrit} : l'orthocentre $H$ d'un triangle est le point d'intersection de ses trois hauteurs ; le centre $O$ de son cercle circonscrit est le point d'intersection des médiatrices de ses côtés. Si $O$ est pris pour origine, $H=A+B+C$.

\section*{Exercices}

\begin{exercise}{Les milieux d'un quadrilatère}
Soit $ABCD$ un quadrilatère du plan ou de l'espace. Soient $P, Q, R, S$ les milieux respectifs de $[AB]$, $[BC]$, $[CD]$, $[DA]$. Montrer que $PQRS$ est un parallélogramme (éventuellement aplati).

\textit{Indication} : Utilisez les expressions vectorielles des milieux et comparez les vecteurs opposés.
\end{exercise}

\begin{exercise}{Hexagone formé par les centroïdes}
Soit \( ABCDEF \) un hexagone quelconque, et soient \( A_1, B_1, C_1, D_1, E_1, F_1 \) les centroïdes des triangles \( ABC, BCD, CDE, DEF, EFA, FAB \). Montrer que les côtés opposés de l’hexagone \( A_1B_1C_1D_1E_1F_1 \) ($[A_1B_1]$ et $[D_1E_1]$, $[B_1C_1]$ et $[E_1F_1]$, $[C_1D_1]$ et $[F_1A_1]$) sont parallèles et de même longueur.

\textit{Indication} : Représentez les centroïdes de chaque triangle et utilisez les propriétés des vecteurs et des translations.
\end{exercise}

\begin{exercise}{Centroïde de $n$ points}
Soient \( A_1, \ldots, A_n \) des points et \( O \) un point quelconque. Montrer qu'il existe un unique point \( S \) tel que
\[
\sum_{i=1}^{n} \vect{SA_i} = \vect{0},
\]
et exprimer \( \vect{OS} \) en fonction des vecteurs \( \vect{OA_i} \).

\textit{Indication} : Exprimez chaque vecteur \( \vect{SA_i} \) en fonction de \( S \) et des points \( A_i \), puis simplifiez la somme vectorielle.
\end{exercise}

\begin{exercise}{Diagonales orthogonales}
Soit $ABCD$ un quadrilatère du plan ou de l'espace. Montrer que ses diagonales $[AC]$ et $[BD]$ sont orthogonales si et seulement si ses deux médianes ont la même longueur.
\end{exercise}

\begin{exercise}{Relation vectorielle dans l'espace}
Soient \( A, B, C, D \) quatre points de l’espace. Démontrer que :
\[
|AB|^2 + |CD|^2 - |BC|^2 - |AD|^2 = 2 \vect{AC} \cdot \vect{DB}.
\]
\end{exercise}

% --- NOUVEAUX EXERCICES AVANCÉS ---
\begin{exercise}{}
Soit $ABC$ un triangle dont tous les angles sont aigus et qui n'est pas équilatéral. Soient $H$ son orthocentre et $O$ le centre de son cercle circonscrit. Montrer que les points $B, H, O, C$ sont sur un même cercle si et seulement si $\angle BAC = 60^\circ$.
\end{exercise}

\begin{exercise}{}
Soit $ABC$ un triangle de centre de gravité $G$. Montrer que
\[
AB^2 + BC^2 + CA^2 = 3\left(GA^2 + GB^2 + GC^2\right).
\]
\end{exercise}

\begin{exercise}{}
Soit $ABC$ un triangle isocèle avec $AB = AC$, de cercle circonscrit de centre $O$, et soit $M$ le milieu de $[BC]$. Montrer que la bissectrice de l'angle $\angle BAC$ est la droite $(AM)$, qu'elle est perpendiculaire à $(BC)$ et qu'elle passe par $O$. En déduire que $O$ est son propre symétrique par rapport à cette bissectrice, et que $O=M$ si et seulement si l'angle en $A$ est droit.
\end{exercise}

\section*{Problèmes}

\begin{problem}{Relation sur les diagonales}
Soient $A, B, C, D$ quatre points du plan ou de l'espace, identifiés à des vecteurs d'origine fixée. Montrer que
\[
|AC|^2 + |BD|^2 = |AB|^2 + |BC|^2 + |CD|^2 + |DA|^2 \Longleftrightarrow A + C = B + D.
\]
(La condition $A+C=B+D$ signifie que $ABCD$ est un parallélogramme, éventuellement aplati.)
\end{problem}

\begin{problem}{Théorème (rectangle dans l’espace)}
Soient \( A, B, C, D \) quatre points dans l’espace. Prouvez le théorème suivant :\\
\textit{Si, pour tout point \( X \) de l’espace, on a}
\[
|AX|^2 + |CX|^2 = |BX|^2 + |DX|^2,
\]
\textit{alors \( ABCD \) est un rectangle} (éventuellement aplati ou réduit à un point).
\end{problem}

\begin{problem}{Théorème d’Euler}
Soit $ABCD$ un quadrilatère du plan ou de l'espace, de médianes $[MN]$ et $[PQ]$ ($M, N, P, Q$ milieux de $[AB], [CD], [BC], [DA]$). Prouvez que
\[
|AC|^2 + |BD|^2 = 2 \left( |MN|^2 + |PQ|^2 \right).
\]
\end{problem}

\begin{problem}{Produit scalaire dans un rectangle}
\textit{Si \( M \) est un point quelconque du plan ou de l'espace et \( ABCD \) un rectangle, alors}
\[
\vect{MA} \cdot \vect{MC} = \vect{MB} \cdot \vect{MD}.
\]
\end{problem}

\begin{problem}{Somme de vecteurs bornés}
Soient \( n \geq 1 \) vecteurs \emph{du plan} \( \vect{a}_1, \ldots, \vect{a}_n \) tels que \( |\vect{a}_i| \leq 1 \). Montrer que, dans la somme
\[
\vect{c} = \pm \vect{a}_1 \pm \cdots \pm \vect{a}_n,
\]
on peut choisir les signes de manière à ce que \( |\vect{c}| \leq \sqrt{2} \).
\end{problem}

\begin{problem}{Triangles équilatéraux et centre de gravité}
Les triangles \( OAB \) et \( OA'B' \) sont équilatéraux et de même orientation. \( S \) est le centre de gravité du triangle \( \triangle OAB \), et \( M \) et \( N \) sont les milieux respectifs de \( A'B \) et \( AB' \). Montrer que
\[
\triangle SMB' \sim \triangle SNA' \quad (\text{IMO Jury 1977}).
\]
(Triangles semblables, de sens opposés : $\dfrac{SM}{SN}=\dfrac{SB'}{SA'}$ et $\widehat{MSB'}=\widehat{NSA'}$.)
\end{problem}

% --- NOUVEAUX PROBLÈMES AVANCÉS ---
\begin{problem}{}
Dans le triangle acutangle $ABC$, soient $O$ le centre du cercle circonscrit et $H$ l'orthocentre. Montrer que la bissectrice intérieure de l'angle $\angle BAC$ est aussi la bissectrice de l'angle $\angle HAO$ (si $H\neq O$, c'est-à-dire si $ABC$ n'est pas équilatéral).
\end{problem}

\begin{problem}{}
Soit $ABC$ un triangle avec $\angle ABC = 2 \angle ACB$ et $\angle BAC < \angle ACB$. Sur le segment $[AB]$, on choisit le point $M$ tel que $\angle MCB = \angle BAC$. Démontrer que la droite $(BC)$ est tangente en $C$ au cercle circonscrit du triangle $AMC$.
\end{problem}

\begin{problem}{}
Un quadrilatère convexe $ABCD$ est donné. Le point $P$ est l'intersection des diagonales $[AC]$ et $[BD]$. Les cercles circonscrits des triangles $ABP$ et $CDP$ se coupent en $P$ et en un second point $Q\neq P$. Prouvez que les triangles $QAB$ et $QCD$ sont semblables.
\end{problem}
